% !TEX program = xelatex
\documentclass{英语试卷}
\begin{document}
\examtitle{2025 届武汉市高三四月调研考试}{英语}
\par 英语学科 试题卷\par
\begin{notice}
\par 1. 答卷前, 考生务必将自己的姓名、准考证号填写在答题卡上。\par
\par 2. 回答选择题时, 选出每小题答案后, 用铅笔把答题卡上对应题目的答案标号涂黑。如需改动, 用橡皮擦干净后, 再选涂其他答案标号。回答非选择题时, 将答案写在答题卡上, 写在本试卷上无效。\par
\par 3. 考试结束后，将本试卷和答题卡一并交回。\par
\end{notice}
\parttitle{第一部分 听力（共两节，满分30分）}
\par 做题时，先将答案标在试卷上，录音内容结束后，你将有两分钟的时间将试卷上的答案转涂到答题卡上。\par
\sectitle{第一节（共5小题：每小题1.5分，满分7.5分）}
\par 听下面 5 段对话。每段对话后有一个小题，从题中所给的 A、B、C 三个选项中选出最佳选项。听完每段对话后，你都有 10 秒钟的时间来回答有关小题和阅读下一小题。每段对话仅读一遍。\par
\question{1. What does the woman advise the man to do?}
\begin{choices}(3)
  \item Continue his task.
  \item Charge his phone.
  \item Have a rest.
\end{choices}
\question{2. What is the man doing?}
\begin{choices}(3)
  \item Playing a video game.
  \item Having a walk outside.
  \item Reading an adventure story.
\end{choices}
\question{3. What are the speakers talking about?}
\begin{choices}(3)
  \item A river.
  \item A film.
  \item A musician.
\end{choices}
\question{4. What does the man have difficulty in?}
\begin{choices}(3)
  \item Raising questions in class.
  \item Keeping up with the material.
  \item Communicating with the teacher.
\end{choices}
\question{5. What does the woman try to master?}
\begin{choices}(3)
  \item Describing feelings.
  \item Combining colours.
  \item Painting landscapes.
\end{choices}
\sectitle{第二节（共15小题：每小题1.5分，满分22.5分）}
\par 听下面 5 段对话或独白。每段对话或独白后有几个小题，从题中所给的 A、B、C 三个选项中选出最佳选项。听每段对话或独白前，你将有时间阅读各个小题，每小题 5 秒钟；听完后，各小题将给出 5 秒钟的作答时间。每段对话或独白读两遍。\par
\par 听第 6 段材料，回答第 6、7 题。\par
\question{6. Where are the speakers?}
\begin{choices}(3)
  \item On a plane.
  \item At a station.
  \item At an airport.
\end{choices}
\question{7. When will the man leave for Boston?}
\begin{choices}(3)
  \item 9:15.
  \item 9:45.
  \item 10:17.
\end{choices}
\par 听第 7 段材料，回答第 8 至 10 题。\par
\question{8. Who is a teacher?}
\begin{choices}(3)
  \item Jenna.
  \item Rachel.
  \item Simon.
\end{choices}
\question{9. How did Jenna become famous?}
\begin{choices}(3)
  \item By launching a blog.
  \item By creating a lifestyle.
  \item By running a magazine.
\end{choices}
\question{10. What will the woman do next?}
\begin{choices}(3)
  \item Visit a friend.
  \item Send a message.
  \item Read a post.
\end{choices}
\par 听第8段材料，回答第11至13题。\par
\question{11. What is the probable relationship between the speakers?}
\begin{choices}(3)
  \item Classmates.
  \item Fellow workers.
  \item Teacher and student.
\end{choices}
\question{12. How does the man feel about Al use at first?}
\begin{choices}(3)
  \item Slightly doubtful.
  \item Rather disappointed.
  \item Pretty satisfied.
\end{choices}
\question{13. What does Al help the woman to do?}
\begin{choices}(3)
  \item Write drafts.
  \item Grade essays.
  \item Generate ideas.
\end{choices}
\par 听第9段材料，回答第14至17题。\par
\question{14. Why did Williams quit football?}
\begin{choices}(3)
  \item He preferred acting.
  \item He hated being pushed.
  \item He was poor at running.
\end{choices}
\question{15. What motivated Williams to start his career in Barcelona?}
\begin{choices}(3)
  \item Gaudi's works.
  \item Local dishes.
  \item Exciting matches.
\end{choices}
\question{16. What challenged Williams the most in Barcelona?}
\begin{choices}(3)
  \item Spanish idioms.
  \item Extreme heat.
  \item Huge crowds.
\end{choices}
\question{17. Which best describe Williams?}
\begin{choices}(3)
  \item Ambitious.
  \item Considerate.
  \item Generous.
\end{choices}
\par 听第 10 段材料，回答第 18 至 20 题。\par
\question{18. What caused the Silver Fire to worsen in California?}
\begin{choices}(3)
  \item Strong winds.
  \item Geographic features.
  \item Human activities.
\end{choices}
\question{19. How many acres has the Table Rock fire burned in South Carolina?}
\begin{choices}(3)
  \item Nearly 574.
  \item Up to 2,000.
  \item About 11,000.
\end{choices}
\question{20. What should certain people in Greenville do?}
\begin{choices}(3)
  \item Fight against hurricanes.
  \item Clean up dry plants.
  \item Keep away from homes.
\end{choices}
\parttitle{第二部分 阅读理解（共两节，满分50分）}
\sectitle{第一节（共15小题：每小题2.5分，满分37.5分）}
\par 阅读下列短文，从每题所给的A、B、C和D四个选项中，选出最佳答案。\par
\passagelabel{A}
\par Terms of Sale\par
\par Subscription Types\par
\par We offer three different types of subscriptions, detailed below.\par
\par\vspace{0.4em}
\begin{center}\begin{tabular}{{|l|p{\dimexpr\linewidth-3em-4\tabcolsep}|}}
\hline
Subscription Types & Provide \\
Digital & · unlimited access to TheAtlantic.com and our iOS and Android app \\
Print \&amp; Digital & · all the benefits of a Digital subscriptions
· 10 print issues of the magazine \\
Premium & · all the benefits of Print \&amp; Digital subscriptions
· ad-free web browsing (浏览) on TheAtlantic.com \\
\hline
\end{tabular}\end{center}
\vspace{0.4em}
\par Subscription Billing Terms\par
\par One-Time Purchases: When you purchase a one-time, non-renewing subscription, we will charge your selected method of payment at the time of purchase.\par
\par Automatically Renewing Subscriptions: Automatically renewing subscriptions require ongoing charges at set times: first, when you initially subscribe and then annually or every two years depending on the option you choose at checkout. You will continue to be charged at your selected interval (间隔) unless, prior to your renewal date, you opt out of automatic renewal or cancel your subscription.\par
\par Gift Subscriptions: When giving a gift subscription, we will charge your selected method of payment at the time of purchase.\par
\par Third-Party Subscriptions: Subscriptions purchased through a third party are only managed through that third party. We are not able to initiate, cancel, or refund any third-party subscription on the customer's behalf. We reserve the right to change our subscription prices at any time.\par
\question{21. What feature is unique to Premium subscriptions?}
\begin{choices}(2)
  \item Limitless website access.
  \item Free iOS and Android app.
  \item Uninterrupted web experiences.
  \item 10 print issues of the magazine.
\end{choices}
\question{22. Which subscription is the magazine unable to cancel directly?}
\begin{choices}(2)
  \item One-Time Purchases.
  \item Automatically Renewing Subscriptions.
  \item Gift Subscriptions.
  \item Third-Party Subscriptions.
\end{choices}
\question{23. What is the purpose of the text?}
\begin{choices}(2)
  \item To persuade.
  \item To inform.
  \item To comment.
  \item To advertise.
\end{choices}
\passagelabel{B}
\par These days, I have been gathering contributions for Class Notes from my classmates at the Brearley School. Seven contributions have come in. They are long or short, emotional or matter-of-fact, and now my job will be to cut and select, rephrase and quote, and count words from each contributor—our entire collection will have to add up to no more than six hundred and fifty words.\par
\par The tough part is trying to give more or less equal space to each of these life accounts, since some of us tend to relate incidents in great detail while others are tight-lipped. If nine of us altogether, including me, contribute notes, then each note is allowed a little over seventy words. If another contribution comes in at the last minute, we all go down to about sixty-five. Do I cut out the trip to Costa Rica, or the visit to the nephew in New Jersey? The recent grandchild, broken leg, or Scrabble (拼字游戏) competition?\par
\par Why did I decide I would not mind being a class-notes agent? I suppose I pursue any piece of writing for two main reasons. First, I like the challenge of trying a new form. In this case, my task is to rearrange the writings of a range of personalities to create one fairly smooth whole. The second reason is always emotional. Here, one source of my emotion was the years-long experience of the school itself. I was also touched that my classmates suggested me for the job, their sensible reason being: "You're a writer."\par
\par As a 65-year-old lady, I feel tired sometimes, but I've never stopped valuing every moment in this job. For me, my classmates are still full of youthful energy, passion, or disgust or tearful hurt. I'm not sure I'm naturally gifted at turning all these colorful details into a brief and entertaining account. That's how I do it—quote their own words as often as I can.\par
\question{24. What does the author need to do with the contributions?}
\begin{choices}(2)
  \item Edit and adjust them.
  \item Collect and grade them.
  \item Ensure the variety of them.
  \item Add personal comments to them.
\end{choices}
\question{25. How does the author explain the difficulty of the job?}
\begin{choices}(2)
  \item By citing quotes.
  \item By listing examples.
  \item By making comparisons.
  \item By giving definitions.
\end{choices}
\question{26. What can we infer about the author from paragraph 3?}
\begin{choices}(2)
  \item She prioritizes new forms.
  \item She prefers emotional writing.
  \item She is hesitant about the role.
  \item She is acknowledged by her classmates.
\end{choices}
\question{27. What does the author think of the job?}
\begin{choices}(2)
  \item It's stressful.
  \item It's worthwhile.
  \item It's straightforward.
  \item It's troublesome.
\end{choices}
\passagelabel{C}
\par Heidi Tait, co-founder of Tangaroa Blue, an organization committed to beach cleaning, was puzzled when she and her teammates started finding strange tire-shaped capsules washed ashore along the Cape York coast, near Australia's northeastern tip. But Tait eventually connected with Satlink—a Spanish satellite communications company. Satlink's GPS-enabled buoys (浮标), the ones the beach cleaners kept finding, help commercial fishers track their nets, lines etc. Tait had a brainwave that would see the buoys not just recycled but reused.\par
\par For more than a decade, boat crews working farther west, in Australia's Gulf of Carpentaria, had been telling Tait about how abandoned fishing nets were circling the gulf, hurting sea creatures. These so-called ghost (幽灵) nets had either broken free from commercial fishing boats and gotten lost, or were cut loose by fishers after getting caught on rocks. Weighing a few tonnes each, the nets that boat crews had chanced upon in the gulf were often too big for them to pull out of the water. They'd typically report the finds to the authorities, but by the time anyone with an appropriately equipped boat could head out to get one, the mass of messy rope had often gone from sight.\par
\par However, once one of the GPS-enabled buoys is attached, it sends its location every few hours to a website, which lets Tangaroa Blue monitor the nets' movements and warns the organization of danger. So far, the reused GPS trackers have helped Tangaroa Blue recover three ghost nets from Australian waters.\par
\par Tangaroa Blue's idea to track ghost nets isn't entirely new; researchers with the Commonwealth Scientific and Industrial Research Organisation (CSIRO), Australia's national science agency, suggested it a decade ago. Taking the idea on board, the authorities tried using battery-powered trackers to monitor a few ghost nets. Tangaroa Blue's approach differs in that its recycled buoys are solar powered and require less upkeep.\par
\question{28. What were the tire-shaped capsules originally used for?}
\begin{choices}(2)
  \item Removing beach waste.
  \item Recycling vehicle tires.
  \item Locating fishing equipment.
  \item Strengthening satellite signals.
\end{choices}
\question{39. What does the author intend to do in paragraph 2?}
\begin{choices}(2)
  \item Illustrate a future possibility.
  \item Describe a technical process.
  \item Present a practical solution.
  \item Frame a problem context.
\end{choices}
\question{30. How do Tangaroa Blue's buoys outperform earlier trackers?}
\begin{choices}(2)
  \item They need less maintenance.
  \item They monitor broader regions.
  \item They receive government funding.
  \item They provide precise tracking data.
\end{choices}
\question{31. What is the best title for the text?}
\begin{choices}(2)
  \item Turning Waste into Treasure: A Fix to Ghost Nets
  \item Fighting Threats: The Hidden Harm of Ghost Nets
  \item Financing Change: Innovative Ocean Cleanup Plans
  \item Protecting Seas from Waste: An Activist of Conservation
\end{choices}
\passagelabel{D}
\par In health research, determining true causes—like whether drinking wine affects health—is vital for informed decisions. But "correlation does not equal causation (因果关系)” is widely recognized. Then how do researchers\par
\par move beyond correlations to find answers?\par
\par Correlations often exist without implying causation. For instance, chocolate consumption correlates with the number of Nobel Prize winners in some countries. Does eating chocolate boost Nobel wins? No. This correlation likely exists because chocolate consumption serves as a proxy for wealth. In turn, wealth relates to education and research funding—key factors behind Nobel achievements. So, just finding a link isn't enough. Scientists need more strong evidence to establish causation.\par
\par Unlike highly controlled chemistry or physics experiments, which can easily show how X affects Y, human biology often relies on indirect evidence. It requires "inductive reasoning"—drawing general conclusions from available evidence. It's like how a prosecutor (检察官) builds a criminal case using circumstantial evidence. While individual pieces might not be persuasive, their build-up strengthens the case. There's one interesting contrast, however. In criminal cases, the standard of proof is "beyond reasonable doubt", but science demands proof based "on the balance of probabilities". This reflects scientists' willingness to revise their beliefs when better evidence emerges.\par
\par Indirect evidence is crucial for inferring causation, but it is not always conclusive. The gold standard for direct evidence of causation is the randomized controlled trial (RCT), where participants are randomly assigned to either receive an intervention or to be a "control". This ensures if you see a difference between the two groups, this can only be due to the effect of the intervention, which effectively proves causation. Unfortunately, ethical (道德的) considerations often prevent us conducting RCTs. For instance, we lack RCT evidence that smoking causes lung cancer because the indirect evidence is so strong that such studies would be unethical.\par
\par Causation in health doesn't work in a simple way. Factors like lifestyle, genes, and environment interact to determine disease risk. Achieving perfect health requires a variety of actions, as no single habit, superfood, or magic pill is the answer.\par
\question{32. What does the underlined word "proxy" in paragraph 2 mean?}
\begin{choices}(2)
  \item Replacement.
  \item Condition.
  \item Indicator.
  \item Preference.
\end{choices}
\question{33. How do scientists' inferences differ from prosecutors' conclusions?}
\begin{choices}(2)
  \item They remain open to update.
  \item They guarantee absolute certainty.
  \item They rely largely on indirect evidence.
  \item They require higher standard of proof.
\end{choices}
\question{34. Why is RTC considered the gold standard for establishing causation?}
\begin{choices}(2)
  \item It offers wide application.
  \item It ensures fair assignment.
  \item It avoids ethical considerations.
  \item It identifies intervention effects.
\end{choices}
\question{35. What is the author's suggestion for health improvement?}
\begin{choices}(2)
  \item Counting on good habits.
  \item Integrating multiple strategies.
  \item Taking a cure-all pill.
  \item Leading an active lifestyle.
\end{choices}
\sectitle{第二节（共5小题；每小题2.5分，满分12.5分）}
\par 阅读下面短文，从短文后的选项中选出能填入空白处的最佳选项。选项中有两项为多余选项。\par
\par I've always disliked getting up in winter. As a child, I would drag myself out of bed, my eyes half-closed, and watch my school uniform warming on the radiator. \blanknum{36} It's too dark, and I am always half-asleep.\par
\par \blanknum{37} Most of them are fun, but few really change my life in the long run. Sunrise alarm clocks, however, are different. They've become a fundamental part of my morning routine. My top picks, the Lumie Bodyclock Spark 100 and Philips SmartSleep, have greatly benefited me. I'm waking up every day feeling more energetic and mentally sharper.\par
\par Essentially, sunrise alarm clocks are a combination of alarm clock and light treatment device that shines with increasing brightness as your wake-up time approaches. The light interacts with our natural sleep-wake cycle, much as the sun does when it rises. 38\par
\par Sunrise alarm clocks also promote healthier sleep habits. In our tech-driven world, waking up often means immediately checking our phones. An important strength of sunrise alarm clocks is that your morning wake-up no longer involves your phone. Stopping a smartphone alarm means your day begins with eyes on-screen. \blanknum{39} But with these clocks, you can start the day with a clearer mind and more concentration.\par
\par Sleep is personal, and no single device fits everyone. But for me, sunrise alarm clocks have been a game-changer. They're affordable, with quality models available for under £30. \blanknum{40} After all, some morning people are made, not born.\par
\begin{choices}(1)
  \item They prove to fall into this unique category.
  \item Interested in tech, I often test fascinating devices.
  \item Therefore, we wake up biologically prepared for the day.
  \item The path to endless screen time is just a fingerprint away.
  \item These days, I set multiple alarms to ensure I'm woken up.
  \item If you are struggling with dark mornings, it's worth a shot.
  \item All of the sunrise alarm clocks I've tested have their strengths.
\end{choices}
\parttitle{第三部分 语言知识运用（共两节，满分30分）}
\sectitle{第一节（共15小题；每小题1分，满分15分）}
\par 阅读下面短文，从短文后各题所给的 A、B、C 和 D 四个选项中，选出可以填入空白处的最佳选项。\par
\par I was excited about the London Olympics and wanted to be part of it. I didn't \blanknum{41} if I was outside the venues dancing. Searching online, I thought they might need \blanknum{42}. Figuring there'd be thousands of applicants for the role, I never expected to \blanknum{43}. Yet weeks later, I was invited to a two-day \blanknum{44} at the London Swimming Centre. It was \blanknum{45}—there were knockout stages and you only got one chance. In the first round, the ones \blanknum{46} to make it to the bottom of a 5-metre-deep pool were cut. Followed were skill exams and theory tests. When I finally passed, I felt \blanknum{47}.\par
\par Online, some \blanknum{48} us as having "the easiest job in the world," but Olympians valued us and sometimes even \blanknum{49} us. I once got a Team GB swimming cap. You could be a top surgeon, but if you don't know how to pull somebody out of the water \blanknum{50}, you've done the damage before they even get treatment. And \blanknum{51} did happen. I remember helping an athlete swimming headfirst into the wall.\par
\par There are starstruck moments. That said, it's very \blanknum{52} when you're in the chair. You forget all about those star athletes because you're so focused on making \blanknum{53} in your head while checking each athlete's movement to see if anything doesn't look right.\par
\par Not everything is shown on TV. Guarding pools is the stuff going on \blanknum{54}. But our silent readiness represents the Games' true spirit, where every role, however \blanknum{55}, matters.\par
\cz{41}{care}{admit}{doubt}{know}
\cz{42}{coaches}{athletes}{lifeguards}{firefighters}
\cz{43}{step away}{hear back}{show off}{sign up}
\cz{44}{routine}{practice}{experiment}{trial}
\cz{45}{dangerous}{simple}{intense}{valuable}
\cz{46}{pretending}{intending}{failing}{desiring}
\cz{47}{misfortune}{disbelief}{dishonesty}{misunderstanding}
\cz{48}{trusted}{admired}{discovered}{dismissed}
\cz{49}{gifted}{assisted}{bothered}{ignored}
\cz{50}{correctly}{immediately}{gently}{hurriedly}
\cz{51}{wonders}{accidents}{changes}{events}
\cz{52}{annoying}{painful}{impressive}{professional}
\cz{53}{health warnings}{risk assessments}{safety rules}{work arrangements}
\cz{54}{behind the scene}{at random}{by chance}{in the spotlight}
\cz{55}{unexpected}{undisturbed}{unhurt}{unseen}
\sectitle{第二节（共10小题；每小题1.5ladnepen分，满分15分）}
\par 阅读下面短文，在空白处填入1个适当的单词或括号内单词的正确形式。\par
\par Renli village in Ziyang city, Sichuan province, Southwest China, is home to a digital nomad community founded by Ke Yu, a girl born in 2000. \blanknum{56} village, just a half-hour drive from downtown, helps escape the pollution, traffic and city crowds. With a gentle mountain wind, the smell of freshly cooked coffee and newly baked bread, it is not hard to figure out \blanknum{57} so many digital nomads have chosen to live and work here.\par
\par "Digital nomads" are generally individuals who, thanks to advances in communication infrastructure (设施), make income through the internet independent of traditional office \blanknum{58} (setting). Ke is familiar with the trend.\par
\par In late 2023, while searching for a remote working location, Ke \blanknum{59} (encounter) a photo online of Renli village with three unused buildings, which, with minimal adaptations, could \blanknum{60} (repurpose) for offices, accommodations and even entertainment spaces for the community.\par
\par Many digital nomads are united by a shared desire, \blanknum{61} is to escape the fast-paced and high-pressure life of big cities to work in scenic and cost-effective locations. But, besides this, what unique experiences can be offered? Ke had an idea.\par
\par Ke and her partners invited artists from around the world \blanknum{62} (design) 10 interconnected art pieces, providing visitors with a \blanknum{63} (true) distinctive artistic journey. Within the 2,000-square-meter shared office space in the community, facilities \blanknum{64} three-dimensional printers, livestreaming rooms and studio lights are provided to meet the specific needs of the remote workers typically \blanknum{65} (pursue) careers such as creative design, programming and media.\par
\parttitle{第四部分 写作（共两节，满分 40 分）}
\sectitle{第一节（满分15分）}
\par 上周你校开展了主题为“科学管理体重”的系列活动，请你给校英文报写一篇报道。内容包括：\par
\par (1) 活动内容;\par
\par (2) 内容反响。\par
\par (1) 写作词数应为 80 左右;\par
\par (2) 请按如下格式在答题纸的相应位置作答。\par
\sectitle{第二节（满分25分）}
\par 阅读下面材料，根据其内容和所给段落开头语续写两段，使之构成一篇完整的短文。\par
\par I hurried into Lehman's grocery store. I was there to buy my grandmother's hearing aid batteries. More importantly, this was my chance to visit the love of my life—Artie.\par
\par I ran to the shelf he was on, and there he was, right where I last saw him: a grey little guy with a large mouth that fit perfectly over my shoulder as I hugged him tightly. I adored that stuffed dolphin. As I held Artie, Mrs. Lehman, who had always been kind to me, passed by and smiled knowingly, "Hugging that dolphin again, huh?"\par
\par But my brother, Rick, noticed I took a long time in Lehman's, saying, "You are a thirteen-year-old girl now.\par
\par You are supposed to be quicker. Next time, I will go with you." Oh, great! My brother, Rick. He always laughed at me.\par
\par But I wouldn't give up. So, as we walked into the store, I suggested to Rick that he get the batteries, since he was older. Then I headed towards Artie. There he was, awaiting me. I grabbed him and hugged him tightly. Just then, I heard Rick saying, "What in the world are you doing?"\par
\par "This is Artie," I said. "I hug him every single time I come here. I can't afford it so I hope he doesn't get sold because I just love this little guy." The look on Rick's face was unreadable, but I knew deep down that he wasn't going to make fun of me. In fact, a few times when he accompanied me to Lehman's, he said, "go on and hug your dolphin, if it's still there."\par
\par One day, Rick was not home, so I went to Lehman's alone. I went straight to Artie's shelf. But he was gone. I searched every shelf but couldn't find him. Suddenly, a voice from the counter interrupted my search. 注意:\par
\par (1) 续写词数应为 150 左右;\par
\par (2) 请按如下格式在答题卡的相应位置作答。\par
\writingpara{Paragraph 1: It was Mrs. Lehman herself.}
\writinglines{8}
\writingpara{Paragraph 2: To my surprise, I found a box lying on my bed.}
\writinglines{8}
\end{document}
